%==============================================================================
% GMAO Application - Documentation Technique Complète
% Guide pour la Compréhension de l'Architecture et du Code
% Date: 31 Janvier 2026
%==============================================================================

\documentclass[12pt,a4paper]{article}

%------------------------------------------------------------------------------
% PACKAGES
%------------------------------------------------------------------------------
\usepackage[utf8]{inputenc}
\usepackage[T1]{fontenc}
\usepackage[french]{babel}
\usepackage{geometry}
\usepackage{graphicx}
\usepackage{xcolor}
\usepackage{tikz}
\usepackage{tcolorbox}
\usepackage{listings}
\usepackage{booktabs}
\usepackage{array}
\usepackage{tabularx}
\usepackage{colortbl}
\usepackage{hyperref}
\usepackage{fontawesome5}
\usepackage{enumitem}
\usepackage{fancyhdr}
\usepackage{titlesec}
\usepackage{float}
\usepackage{caption}
\usepackage{subcaption}
\usepackage{amsmath}
\usepackage{amssymb}
\usepackage{multicol}

%------------------------------------------------------------------------------
% CONFIGURATION DE LA PAGE
%------------------------------------------------------------------------------
\geometry{
    left=2cm,
    right=2cm,
    top=2.5cm,
    bottom=2.5cm
}

%------------------------------------------------------------------------------
% COULEURS PERSONNALISÉES
%------------------------------------------------------------------------------
% Couleurs principales
\definecolor{primaryblue}{RGB}{41, 98, 255}
\definecolor{secondarygreen}{RGB}{46, 204, 113}
\definecolor{accentorange}{RGB}{243, 156, 18}
\definecolor{warningred}{RGB}{231, 76, 60}
\definecolor{darkgray}{RGB}{52, 73, 94}
\definecolor{lightgray}{RGB}{236, 240, 241}
\definecolor{codebg}{RGB}{45, 45, 45}
\definecolor{codegreen}{RGB}{152, 195, 121}
\definecolor{codepurple}{RGB}{198, 120, 221}
\definecolor{codeyellow}{RGB}{229, 192, 123}
\definecolor{codeblue}{RGB}{97, 175, 239}

% Couleurs spécifiques aux technologies
\definecolor{dotnetpurple}{RGB}{81, 43, 212}
\definecolor{angularred}{RGB}{221, 0, 49}
\definecolor{pythonyellow}{RGB}{255, 212, 59}
\definecolor{dockerblue}{RGB}{33, 150, 243}

% Couleurs pour les boîtes
\definecolor{infobox}{RGB}{52, 152, 219}
\definecolor{successbox}{RGB}{39, 174, 96}
\definecolor{warningbox}{RGB}{241, 196, 15}
\definecolor{dangerbox}{RGB}{192, 57, 43}

%------------------------------------------------------------------------------
% CONFIGURATION DES LIENS
%------------------------------------------------------------------------------
\hypersetup{
    colorlinks=true,
    linkcolor=primaryblue,
    urlcolor=primaryblue,
    citecolor=primaryblue
}

%------------------------------------------------------------------------------
% STYLE DES TITRES
%------------------------------------------------------------------------------
\titleformat{\section}
    {\color{primaryblue}\LARGE\bfseries}
    {\thesection}{1em}{}[\titlerule]

\titleformat{\subsection}
    {\color{darkgray}\Large\bfseries}
    {\thesubsection}{1em}{}

\titleformat{\subsubsection}
    {\color{darkgray}\large\bfseries}
    {\thesubsubsection}{1em}{}

%------------------------------------------------------------------------------
% BOÎTES COLORÉES (tcolorbox)
%------------------------------------------------------------------------------
\tcbuselibrary{skins,breakable}

% Boîte Info
\newtcolorbox{infobox}[1][]{
    enhanced,
    colback=infobox!10,
    colframe=infobox,
    fonttitle=\bfseries,
    title={\faInfoCircle\ Information},
    #1
}

% Boîte Succès
\newtcolorbox{successbox}[1][]{
    enhanced,
    colback=successbox!10,
    colframe=successbox,
    fonttitle=\bfseries,
    title={\faCheckCircle\ Point Clé},
    #1
}

% Boîte Avertissement
\newtcolorbox{warningbox}[1][]{
    enhanced,
    colback=warningbox!10,
    colframe=warningbox,
    fonttitle=\bfseries,
    title={\faExclamationTriangle\ Important},
    #1
}

% Boîte Définition
\newtcolorbox{defbox}[2][]{
    enhanced,
    colback=lightgray,
    colframe=primaryblue,
    fonttitle=\bfseries,
    title={\faBook\ #2},
    #1
}

% Boîte Code/Technique
\newtcolorbox{techbox}[2][]{
    enhanced,
    colback=codebg,
    colframe=darkgray,
    fonttitle=\bfseries\color{white},
    coltitle=white,
    title={\faCode\ #2},
    #1
}

% Boîte .NET
\newtcolorbox{dotnetbox}[1][]{
    enhanced,
    colback=dotnetpurple!10,
    colframe=dotnetpurple,
    fonttitle=\bfseries,
    title={\faServer\ .NET Core},
    #1
}

% Boîte Angular
\newtcolorbox{angularbox}[1][]{
    enhanced,
    colback=angularred!10,
    colframe=angularred,
    fonttitle=\bfseries,
    title={\faAngular\ Angular},
    #1
}

% Boîte Python
\newtcolorbox{pythonbox}[1][]{
    enhanced,
    colback=pythonyellow!15,
    colframe=accentorange,
    fonttitle=\bfseries,
    title={\faPython\ Python FastAPI},
    #1
}

%------------------------------------------------------------------------------
% CONFIGURATION DES LISTINGS (CODE)
%------------------------------------------------------------------------------
\lstset{
    backgroundcolor=\color{codebg},
    basicstyle=\ttfamily\small\color{white},
    keywordstyle=\color{codepurple}\bfseries,
    stringstyle=\color{codegreen},
    commentstyle=\color{gray}\itshape,
    numberstyle=\tiny\color{gray},
    numbers=left,
    numbersep=8pt,
    breaklines=true,
    breakatwhitespace=true,
    tabsize=4,
    showstringspaces=false,
    frame=single,
    rulecolor=\color{darkgray},
    captionpos=b,
    xleftmargin=15pt,
    framexleftmargin=15pt
}

\lstdefinelanguage{CSharp}{
    keywords={public, private, protected, class, interface, namespace, using, async, await, return, var, new, if, else, for, foreach, while, try, catch, finally, throw, static, readonly, const, override, virtual, abstract, sealed, partial, get, set, int, string, bool, void, Task, ActionResult, IActionResult, HttpGet, HttpPost, HttpPut, HttpDelete, Route, ApiController, Authorize},
    keywordstyle=\color{codeblue}\bfseries,
    sensitive=true,
    morestring=[b]",
    morecomment=[l]{//},
    morecomment=[s]{/*}{*/},
}

\lstdefinelanguage{TypeScript}{
    keywords={import, export, from, class, interface, extends, implements, constructor, public, private, protected, readonly, async, await, return, const, let, var, if, else, for, while, function, new, this, true, false, null, undefined, Observable, BehaviorSubject, Component, Injectable, OnInit, subscribe, pipe, tap, map, catchError},
    keywordstyle=\color{codeblue}\bfseries,
    sensitive=true,
    morestring=[b]",
    morestring=[b]',
    morestring=[b]`,
    morecomment=[l]{//},
    morecomment=[s]{/*}{*/},
}

\lstdefinelanguage{Python}{
    keywords={def, return, if, else, elif, for, while, in, import, from, as, class, try, except, with, True, False, None, and, or, not, async, await, self},
    keywordstyle=\color{codepurple}\bfseries,
    sensitive=true,
    morestring=[b]",
    morestring=[b]',
    morecomment=[l]{\#},
}

\lstdefinelanguage{SQL}{
    keywords={SELECT, FROM, WHERE, JOIN, LEFT, RIGHT, INNER, ON, AND, OR, NOT, IN, AS, TOP, ORDER, BY, GROUP, HAVING, COUNT, SUM, AVG, MAX, MIN, INSERT, UPDATE, DELETE, CREATE, TABLE},
    keywordstyle=\color{codeblue}\bfseries,
    sensitive=false,
    morestring=[b]',
    morecomment=[l]{--},
}

%------------------------------------------------------------------------------
% EN-TÊTE ET PIED DE PAGE
%------------------------------------------------------------------------------
\pagestyle{fancy}
\fancyhf{}
\fancyhead[L]{\textcolor{darkgray}{\small GMAO Application - Documentation Complète}}
\fancyhead[R]{\textcolor{darkgray}{\small \thepage}}
\fancyfoot[C]{\textcolor{gray}{\small Documentation Technique - Janvier 2026}}
\renewcommand{\headrulewidth}{0.4pt}
\renewcommand{\footrulewidth}{0.4pt}

%------------------------------------------------------------------------------
% TikZ LIBRARIES
%------------------------------------------------------------------------------
\usetikzlibrary{shapes.geometric, arrows.meta, positioning, calc, shadows, fit, backgrounds, decorations.pathreplacing}

%==============================================================================
% DÉBUT DU DOCUMENT
%==============================================================================
\begin{document}

%------------------------------------------------------------------------------
% PAGE DE TITRE
%------------------------------------------------------------------------------
\begin{titlepage}
    \centering
    \vspace*{1cm}
    
    % Logo/Icône
    {\Huge\color{primaryblue}\faCogs}
    
    \vspace{0.8cm}
    
    {\Huge\bfseries\color{primaryblue} Application GMAO}
    
    \vspace{0.3cm}
    
    {\LARGE\color{darkgray} Gestion de Maintenance Assistée par Ordinateur}
    
    \vspace{1.5cm}
    
    \begin{tikzpicture}
        \node[draw=primaryblue, fill=primaryblue!10, rounded corners=10pt, 
              minimum width=13cm, minimum height=4cm, align=center] {
            {\Large\bfseries Documentation Technique Complète}\\[0.4cm]
            {\large Architecture $\bullet$ Technologies $\bullet$ Code Source}\\[0.3cm]
            {\large Workflow $\bullet$ Modules $\bullet$ DevOps}
        };
    \end{tikzpicture}
    
    \vspace{1.5cm}
    
    % Technologies utilisées
    \begin{tikzpicture}
        \node[draw=dotnetpurple, fill=dotnetpurple!20, rounded corners=5pt, minimum width=2.5cm, minimum height=1cm] at (0,0) {\textbf{.NET 9}};
        \node[draw=angularred, fill=angularred!20, rounded corners=5pt, minimum width=2.5cm, minimum height=1cm] at (3.5,0) {\textbf{Angular 17}};
        \node[draw=accentorange, fill=pythonyellow!30, rounded corners=5pt, minimum width=2.5cm, minimum height=1cm] at (7,0) {\textbf{FastAPI}};
        \node[draw=darkgray, fill=darkgray!20, rounded corners=5pt, minimum width=2.5cm, minimum height=1cm] at (10.5,0) {\textbf{SQL Server}};
    \end{tikzpicture}
    
    \vspace{1.5cm}
    
    \begin{tabular}{ll}
        \faCalendar & \textbf{Date :} 31 Janvier 2026 \\[0.3cm]
        \faCode & \textbf{Version :} 1.0 \\[0.3cm]
        \faUser & \textbf{Auteur :} Stagiaire GMAO \\[0.3cm]
        \faBuilding & \textbf{Projet :} Stage de Fin d'Études
    \end{tabular}
    
    \vfill
    
    {\large\color{gray} Documentation pour la Soutenance de Stage}
    
\end{titlepage}

%------------------------------------------------------------------------------
% TABLE DES MATIÈRES
%------------------------------------------------------------------------------
\tableofcontents
\newpage

%==============================================================================
% SECTION 1 : VUE D'ENSEMBLE
%==============================================================================
\section{\faEye\ Vue d'Ensemble de l'Application}

\subsection{Qu'est-ce que le GMAO ?}

\begin{defbox}{GMAO - Gestion de Maintenance Assistée par Ordinateur}
Le \textbf{GMAO} est un système informatique permettant de gérer l'ensemble des activités de maintenance d'une entreprise :
\begin{itemize}
    \item \textbf{Gestion des équipements} : inventaire, caractéristiques, historique
    \item \textbf{Gestion des interventions} : demandes, planification, suivi
    \item \textbf{Gestion des stocks} : pièces de rechange, fournisseurs
    \item \textbf{Gestion des ressources} : techniciens, compétences
    \item \textbf{Assistant IA} : chatbot pour interroger les données
\end{itemize}
\end{defbox}

\subsection{Objectifs du Projet}

\begin{successbox}[title={\faCheckCircle\ Objectifs Principaux}]
\begin{enumerate}
    \item Créer une application web moderne pour la gestion de maintenance
    \item Implémenter un système d'authentification sécurisé avec rôles
    \item Développer un chatbot IA capable d'interroger la base de données
    \item Préparer l'infrastructure pour un déploiement DevOps
\end{enumerate}
\end{successbox}

\subsection{Architecture Globale}

\begin{figure}[H]
\centering
\begin{tikzpicture}[
    scale=0.8,
    transform shape,
    node distance=1cm,
    component/.style={rectangle, rounded corners=8pt, draw=#1, fill=#1!20,
                      minimum width=9cm, minimum height=1.3cm, align=center,
                      font=\small, drop shadow},
    arrow/.style={-{Stealth[length=3mm]}, thick, color=darkgray},
    nodelabel/.style={font=\footnotesize, color=gray}
]

% User
\node[rectangle, rounded corners=5pt, draw=primaryblue, fill=primaryblue!30,
      minimum width=3cm, minimum height=0.8cm] (user) at (0, 6) {
    \faUser\ \textbf{Utilisateur}
};

% Frontend
\node[component=angularred] (frontend) at (0, 4) {
    \textcolor{angularred}{\faAngular\ \textbf{FRONTEND - Angular 17}}\\
    SPA (Single Page Application) - Port 4200
};

% API
\node[component=dotnetpurple] (api) at (0, 2) {
    \textcolor{dotnetpurple}{\faServer\ \textbf{API - .NET Core 9}}\\
    RESTful API + Cookie Auth - Port 5000
};

% Python Service
\node[component=accentorange] (python) at (5, 0) {
    \textcolor{accentorange}{\faPython\ \textbf{CHATBOT}}\\
    FastAPI - Port 5002
};

% Database
\node[component=darkgray] (db) at (0, 0) {
    \textcolor{darkgray}{\faDatabase\ \textbf{SQL SERVER}}\\
    Base de données GMAO
};

% LLM
\node[rectangle, rounded corners=5pt, draw=codepurple, fill=codepurple!20,
      minimum width=3cm, minimum height=1cm, align=center] (llm) at (9, 0) {
    \faBrain\ \textbf{LLM} \\ OpenRouter
};

% Arrows
\draw[arrow, <->] (user) -- (frontend);
\draw[arrow, <->] (frontend) -- node[right, nodelabel] {HTTP/JSON} (api);
\draw[arrow, <->] (api) -- node[left, nodelabel] {EF Core} (db);
\draw[arrow, <->] (api) -- node[above, nodelabel] {HTTP} (python);
\draw[arrow, <->] (python) -- node[above, nodelabel] {PyODBC} (db);
\draw[arrow, <->] (python) -- node[above, nodelabel] {API} (llm);

\end{tikzpicture}
\caption{Architecture Multi-Tiers de l'Application GMAO}
\end{figure}

\newpage

%==============================================================================
% SECTION 2 : TECHNOLOGIES & CONCEPTS
%==============================================================================
\section{\faTools\ Technologies et Concepts Clés}

\subsection{Backend : .NET Core 9}

\begin{dotnetbox}
\textbf{.NET Core} est un framework open-source de Microsoft pour développer des applications web modernes.

\begin{itemize}
    \item \textbf{Cross-platform} : Fonctionne sur Windows, Linux, macOS
    \item \textbf{Performance} : Très rapide grâce à la compilation JIT
    \item \textbf{Écosystème riche} : NuGet packages, Entity Framework
\end{itemize}
\end{dotnetbox}

\subsubsection{Entity Framework Core (ORM)}

\begin{defbox}{ORM - Object-Relational Mapping}
\textbf{Entity Framework Core} est un ORM qui permet de manipuler la base de données avec des objets C\# au lieu d'écrire du SQL brut.

\textbf{Avantages :}
\begin{itemize}
    \item Code plus lisible et maintenable
    \item Protection contre les injections SQL
    \item Migrations automatiques de la base de données
\end{itemize}
\end{defbox}

\begin{lstlisting}[language=CSharp, caption=Exemple Entity Framework - Récupération des interventions]
// Sans EF Core (SQL brut) - DANGEREUX et verbeux
string sql = "SELECT * FROM DemandeInterventions WHERE IsActif = 1";

// Avec EF Core (LINQ) - SECURISE et lisible
var interventions = await _context.DemandeInterventions
    .Where(d => d.IsActif)           // Filtre: WHERE IsActif = 1
    .Include(d => d.Status)          // JOIN avec table Status
    .Include(d => d.Demandeur)       // JOIN avec table Users
    .ToListAsync();                  // Execute et retourne List<>
\end{lstlisting}

\begin{infobox}[title={\faLightbulb\ Explication Ligne par Ligne}]
\begin{itemize}
    \item \texttt{\_context} : Instance du DbContext injectée, représente la connexion à la DB
    \item \texttt{.Where(d => d.IsActif)} : Expression Lambda - \texttt{d} est chaque élément, \texttt{=>} signifie "tel que"
    \item \texttt{.Include()} : \textbf{Eager Loading} - charge les relations en une seule requête (évite le problème N+1)
    \item \texttt{await ... ToListAsync()} : Exécution asynchrone - ne bloque pas le thread pendant la requête DB
\end{itemize}
\end{infobox}

\subsubsection{ASP.NET Identity (Authentification)}

\begin{defbox}{Identity - Système d'Authentification}
\textbf{ASP.NET Identity} gère :
\begin{itemize}
    \item La création et gestion des utilisateurs
    \item Le hashage sécurisé des mots de passe
    \item La gestion des rôles (Superviseur, Demandeur, Utilisateur)
    \item Les tokens de réinitialisation de mot de passe
\end{itemize}
\end{defbox}

\subsubsection{Cookie Authentication vs JWT}

\begin{warningbox}[title={\faKey\ Authentification par Cookie (Notre Choix)}]
Notre application utilise l'authentification par \textbf{Cookie} plutôt que JWT :

\textbf{Avantages des Cookies :}
\begin{itemize}
    \item Gérés automatiquement par le navigateur
    \item Plus sécurisés (HttpOnly, Secure flags)
    \item Expiration automatique avec sliding expiration
    \item Pas besoin de stocker le token côté client
\end{itemize}

\textbf{Configuration :}
\begin{itemize}
    \item \texttt{HttpOnly = true} : Cookie inaccessible par JavaScript
    \item \texttt{Secure = true} : Cookie envoyé uniquement en HTTPS
    \item \texttt{SameSite = None} : Pour les requêtes cross-origin
\end{itemize}
\end{warningbox}

\begin{lstlisting}[language=CSharp, caption=Configuration Cookie Authentication]
builder.Services.ConfigureApplicationCookie(options => 
{
    options.Cookie.SecurePolicy = CookieSecurePolicy.Always;
    options.Cookie.SameSite = SameSiteMode.None;
    options.Cookie.HttpOnly = true;
    options.ExpireTimeSpan = TimeSpan.FromDays(7);
    options.SlidingExpiration = true;
});
\end{lstlisting}

\begin{infobox}[title={\faLightbulb\ Explication de Chaque Option}]
\begin{itemize}
    \item \texttt{SecurePolicy.Always} : Cookie transmis \textbf{uniquement en HTTPS} - empêche l'interception sur HTTP non sécurisé
    \item \texttt{SameSite.None} : Permet l'envoi du cookie dans les requêtes \textbf{cross-origin} (Angular sur port 4200 → API sur port 5000)
    \item \texttt{HttpOnly = true} : Cookie \textbf{inaccessible via JavaScript} (\texttt{document.cookie}) - protection contre les attaques XSS
    \item \texttt{ExpireTimeSpan} : Durée de vie du cookie (7 jours) - après, l'utilisateur doit se reconnecter
    \item \texttt{SlidingExpiration = true} : Chaque requête \textbf{renouvelle} l'expiration - l'utilisateur actif reste connecté
\end{itemize}
\end{infobox}

\newpage

\subsection{Frontend : Angular 17+}

\begin{angularbox}
\textbf{Angular} est un framework TypeScript pour créer des SPA (Single Page Applications).

\textbf{Caractéristiques clés :}
\begin{itemize}
    \item \textbf{Component-based} : Application divisée en composants réutilisables
    \item \textbf{TypeScript} : JavaScript avec typage statique
    \item \textbf{Reactive} : Programmation réactive avec RxJS
    \item \textbf{Standalone Components} : Nouvelle syntaxe Angular 17+
\end{itemize}
\end{angularbox}

\subsubsection{Concepts Angular Essentiels}

\begin{defbox}{Observable (RxJS)}
Un \textbf{Observable} représente un flux de données asynchrones. Contrairement à une Promise qui retourne une seule valeur, un Observable peut émettre plusieurs valeurs dans le temps.

\begin{lstlisting}[language=TypeScript]
// Promise : une seule valeur
fetch('/api/data').then(data => console.log(data));

// Observable : flux de valeurs
this.http.get('/api/data').subscribe(data => {
    console.log(data); // Peut etre appele plusieurs fois
});
\end{lstlisting}
\end{defbox}

\begin{defbox}{BehaviorSubject}
Un \textbf{BehaviorSubject} est un type spécial d'Observable qui :
\begin{itemize}
    \item Garde en mémoire la dernière valeur émise
    \item Émet immédiatement cette valeur aux nouveaux abonnés
    \item Peut être mis à jour avec \texttt{.next(value)}
\end{itemize}

\begin{lstlisting}[language=TypeScript]
// Declaration avec valeur initiale
private currentUserSubject = new BehaviorSubject<User | null>(null);
public currentUser$ = this.currentUserSubject.asObservable();

// Mise a jour
this.currentUserSubject.next(newUser);

// Lecture de la valeur actuelle
const user = this.currentUserSubject.value;
\end{lstlisting}
\end{defbox}

\begin{defbox}{Interceptor HTTP}
Un \textbf{Interceptor} est un middleware qui intercepte toutes les requêtes HTTP sortantes et/ou les réponses entrantes.

\textbf{Utilisations :}
\begin{itemize}
    \item Ajouter automatiquement les credentials (cookies)
    \item Gérer les erreurs 401 (redirection vers login)
    \item Ajouter des headers communs
    \item Logging des requêtes
\end{itemize}
\end{defbox}

\begin{lstlisting}[language=TypeScript, caption=Auth Interceptor - Gestion automatique des cookies]
export const authInterceptor: HttpInterceptorFn = (req, next) => {
  // Verifier si la requete va vers notre API
  if (req.url.includes(environment.apiUrl)) {
    // Cloner la requete avec credentials
    req = req.clone({
      withCredentials: true
    });
  }

  // Passer au prochain handler et gerer les erreurs
  return next(req).pipe(
    catchError((error) => {
      if (error.status === 401) {
        const router = inject(Router);
        router.navigate(['/login']);
      }
      return throwError(() => error);
    })
  );
};
\end{lstlisting}

\begin{infobox}[title={\faLightbulb\ Explication Détaillée de l'Interceptor}]
\textbf{Ligne 1 - Signature :}
\begin{itemize}
    \item \texttt{HttpInterceptorFn} : Type fonction d'interceptor (Angular 17+ standalone)
    \item \texttt{req} : La requête HTTP sortante (objet immuable)
    \item \texttt{next} : Fonction pour passer au handler suivant
\end{itemize}

\textbf{Ligne 3-7 - Modification de la requête :}
\begin{itemize}
    \item \texttt{req.url.includes()} : Vérifie si l'URL contient notre API
    \item \texttt{req.clone(\{...\})} : Crée une \textbf{copie} avec modifications (requête immuable)
    \item \texttt{withCredentials: true} : Inclut les cookies dans la requête cross-origin
\end{itemize}

\textbf{Ligne 10-17 - Pipeline RxJS :}
\begin{itemize}
    \item \texttt{next(req)} : Envoie la requête, retourne un Observable
    \item \texttt{.pipe()} : Chaîne d'opérateurs RxJS
    \item \texttt{catchError()} : Intercepte les erreurs HTTP
    \item \texttt{inject(Router)} : Injection de dépendance dans une fonction
    \item \texttt{throwError(() => error)} : Propage l'erreur aux autres handlers
\end{itemize}
\end{infobox}

\begin{defbox}{Data Binding (Liaison de Données)}
Angular offre plusieurs types de data binding :

\begin{itemize}
    \item \textbf{Interpolation} \texttt{\{\{value\}\}} : Affiche une valeur dans le template
    \item \textbf{Property Binding} \texttt{[property]="value"} : Lie une propriété HTML
    \item \textbf{Event Binding} \texttt{(event)="handler()"} : Écoute un événement
    \item \textbf{Two-Way Binding} \texttt{[(ngModel)]="value"} : Liaison bidirectionnelle
\end{itemize}

\begin{lstlisting}[language=TypeScript]
<!-- Interpolation -->
<p>Bienvenue, {{user.name}}</p>

<!-- Property Binding -->
<button [disabled]="isLoading">Envoyer</button>

<!-- Event Binding -->
<button (click)="sendMessage()">Envoyer</button>

<!-- Two-Way Binding -->
<input [(ngModel)]="searchTerm" />
\end{lstlisting}
\end{defbox}

\newpage

\subsection{Chatbot : Python FastAPI}

\begin{pythonbox}
\textbf{FastAPI} est un framework Python moderne pour créer des APIs REST.

\textbf{Caractéristiques :}
\begin{itemize}
    \item \textbf{Très rapide} : Basé sur Starlette et Pydantic
    \item \textbf{Documentation auto} : Swagger UI automatique
    \item \textbf{Validation} : Validation automatique avec Pydantic
    \item \textbf{Async} : Support natif de l'asynchrone
\end{itemize}
\end{pythonbox}

\subsubsection{Pydantic (Validation)}

\begin{defbox}{Pydantic - Validation des Données}
\textbf{Pydantic} valide automatiquement les données entrantes et sortantes.

\begin{lstlisting}[language=Python]
from pydantic import BaseModel
from typing import Optional, List, Dict

class ChatRequest(BaseModel):
    query: str  # Requis
    conversation_history: Optional[List[Dict[str, str]]] = None

class ChatResponse(BaseModel):
    response: str
    metadata: Optional[Dict[str, Any]] = None
    suggestions: Optional[List[str]] = None
\end{lstlisting}
\end{defbox}

\newpage

%==============================================================================
% SECTION 3 : MODULE INTERVENTIONS (DÉTAILLÉ)
%==============================================================================
\section{\faWrench\ Module Interventions - Étude Détaillée}

\subsection{Vue d'Ensemble du Module}

Le module \textbf{Interventions} est le cœur de l'application GMAO. Il permet de gérer tout le cycle de vie d'une demande d'intervention.

\begin{figure}[H]
\centering
\begin{tikzpicture}[
    node distance=0.8cm,
    state/.style={rectangle, rounded corners=8pt, draw=#1, fill=#1!30,
                  minimum width=2.5cm, minimum height=1cm, align=center, font=\small},
    arrow/.style={-{Stealth[length=3mm]}, thick}
]

% States
\node[state=primaryblue] (new) {Nouveau};
\node[state=accentorange, right=2cm of new] (assigned) {Assignée};
\node[state=infobox, right=2cm of assigned] (progress) {En cours};
\node[state=successbox, right=2cm of progress] (done) {Terminée};
\node[state=warningred, below=1.5cm of assigned] (cancel) {Annulée};

% Arrows
\draw[arrow, primaryblue] (new) -- node[above, font=\footnotesize] {Assignation} (assigned);
\draw[arrow, accentorange] (assigned) -- node[above, font=\footnotesize] {Prise en charge} (progress);
\draw[arrow, infobox] (progress) -- node[above, font=\footnotesize] {Résolution} (done);
\draw[arrow, warningred, dashed] (new) -- (cancel);
\draw[arrow, warningred, dashed] (assigned) -- (cancel);

\end{tikzpicture}
\caption{Workflow d'une Demande d'Intervention}
\end{figure}

\subsection{Entité : DemandeIntervention}

\begin{lstlisting}[language=CSharp, caption=Core/Entities/interventions/DemandeIntervention.cs]
namespace Core.Entities
{
    public class DemandeIntervention
    {
        // Cle primaire
        public int IdDemandeIntervention { get; set; }
        
        // Informations principales
        public string Titre { get; set; } = string.Empty;
        public DateTime DateCreation { get; set; }
        public bool MachineArretee { get; set; }
        
        // Fichier attache (stocke en binaire)
        public byte[]? FileData { get; set; }
        public string? FileName { get; set; }
        public string? ContentType { get; set; }

        // Relations (Foreign Keys)
        public int StatusId { get; set; }
        public Status Status { get; set; } = null!;
        
        public int IDTypeIntervention { get; set; }
        public TypeIntervention TypeIntervention { get; set; } = null!;
        
        public int IDUrgence { get; set; }
        public Urgence Urgence { get; set; } = null!;
        
        // Utilisateurs lies
        public int DemandeursID { get; set; }
        public Users Demandeur { get; set; } = null!;
        
        public int? DestinatairesID { get; set; }  // Nullable = optionnel
        public Users? Destinataire { get; set; }
        
        // Suivi du workflow
        public DateTime? DateAssignation { get; set; }
        public DateTime? DatePriseEnCharge { get; set; }
        public DateTime? DateResolution { get; set; }
        public DateTime? DateCloture { get; set; }
        public string? RaisonAnnulation { get; set; }
        
        public bool IsActif { get; set; }  // Soft delete
    }
}
\end{lstlisting}

\begin{infobox}[title={\faLightbulb\ Explication Détaillée du Code}]
\textbf{Propriétés de Base :}
\begin{itemize}
    \item \texttt{public int IdDemandeIntervention \{ get; set; \}} : Propriété auto-implémentée en C\# - le compilateur génère automatiquement le champ privé
    \item \texttt{= string.Empty} : Initialisation par défaut pour éviter les null
\end{itemize}

\textbf{Navigation Properties (Relations) :}
\begin{itemize}
    \item \texttt{public int StatusId} : Clé étrangère (FK) - stocke l'ID de la relation
    \item \texttt{public Status Status} : Navigation Property - permet d'accéder à l'objet complet
    \item \texttt{= null!} : Le \texttt{!} dit au compilateur "je sais que ce n'est pas null" (évite les warnings)
\end{itemize}

\textbf{Types Nullable (?) :}
\begin{itemize}
    \item \texttt{int?} = \texttt{Nullable<int>} : Peut contenir un int OU null
    \item \texttt{DestinatairesID?} : Le technicien assigné est optionnel (pas toujours assigné)
    \item \texttt{DateTime?} : Les dates de workflow sont null tant que l'action n'est pas faite
\end{itemize}

\textbf{Patterns Importants :}
\begin{itemize}
    \item \textbf{Soft Delete} : \texttt{IsActif = false} au lieu de DELETE - conserve l'historique
    \item \textbf{Binary Storage} : \texttt{byte[]?} pour stocker des fichiers directement en DB
\end{itemize}
\end{infobox}

\newpage

\subsection{DTO : Data Transfer Object}

\begin{defbox}{DTO - Data Transfer Object}
Un \textbf{DTO} est un objet utilisé pour transférer des données entre les couches de l'application. Il permet de :
\begin{itemize}
    \item Ne pas exposer directement les entités de la base
    \item Valider les données entrantes
    \item Contrôler les champs envoyés/reçus
\end{itemize}
\end{defbox}

\begin{lstlisting}[language=CSharp, caption=API/DTOs/InterventionDto.cs]
using System.Text.Json.Serialization;

namespace API.DTOs
{
    public class InterventionDto
    {
        public int IdDemandeIntervention { get; set; }
        
        [JsonPropertyName("titre")]  // Nom JSON different du C#
        public string? Titre { get; set; }
        
        [JsonPropertyName("dateCreation")]
        public DateTime DateCreation { get; set; }
        
        [JsonPropertyName("machineArretee")]
        public bool MachineArretee { get; set; }
        
        [JsonPropertyName("idTypeIntervention")]
        public int? IDTypeIntervention { get; set; }
        
        [JsonPropertyName("idUrgence")]
        public int? IDUrgence { get; set; }
        
        [JsonPropertyName("demandeursId")]
        public int? DemandeursID { get; set; }
        
        [JsonPropertyName("destinatairesId")]
        public int? DestinatairesID { get; set; }
        
        // Actions de workflow
        public string? WorkflowAction { get; set; }
        public string? RaisonAnnulation { get; set; }
    }
}
\end{lstlisting}

\begin{infobox}[title={\faLightbulb\ Explication Complète du DTO}]
\textbf{Pourquoi un DTO séparé de l'Entity ?}
\begin{itemize}
    \item \textbf{Sécurité} : Ne pas exposer tous les champs (ex: FileData binaire)
    \item \textbf{Validation} : Contrôler ce que le client peut envoyer
    \item \textbf{Découplage} : Changer l'Entity sans casser l'API
\end{itemize}

\textbf{L'attribut [JsonPropertyName] :}
\begin{itemize}
    \item C\# utilise \texttt{PascalCase} : \texttt{Titre}, \texttt{DateCreation}
    \item JavaScript/JSON utilise \texttt{camelCase} : \texttt{titre}, \texttt{dateCreation}
    \item L'attribut fait la \textbf{correspondance automatique} lors de la sérialisation
\end{itemize}

\textbf{Types Nullable (?) dans le DTO :}
\begin{itemize}
    \item \texttt{string?} : Le champ est optionnel dans la requête
    \item \texttt{int?} : Permet de distinguer "non fourni" (null) de "valeur 0"
    \item Utile pour les \textbf{PATCH} partiels où on ne met à jour que certains champs
\end{itemize}
\end{infobox}

\newpage

\subsection{Controller : API REST}

\begin{lstlisting}[language=CSharp, caption=DemandeInterventionController.cs - Structure]
[Authorize]  // Requiert authentification
[ApiController]
[Route("api/[controller]")]  // Route: /api/demandeintervention
public class DemandeInterventionController : ControllerBase
{
    private readonly StoreContext _context;  // Acces DB
    private readonly UserManager<AppUser> _userManager;  // Gestion users
    private readonly ILogger<DemandeInterventionController> _logger;
    
    // Injection de dependances via constructeur
    public DemandeInterventionController(
        StoreContext context, 
        UserManager<AppUser> userManager,
        ILogger<DemandeInterventionController> logger)
    {
        _context = context;
        _userManager = userManager;
        _logger = logger;
    }
}
\end{lstlisting}

\begin{defbox}{Injection de Dépendances (DI) - Concept Fondamental}
L'\textbf{Injection de Dépendances} est un pattern où les dépendances sont \textbf{fournies} à une classe plutôt que créées par elle-même.

\textbf{Sans DI (Mauvaise pratique) :}
\begin{lstlisting}[language=CSharp, numbers=none]
public class MyController {
    private StoreContext _context = new StoreContext(); // Mauvais!
}
\end{lstlisting}

\textbf{Avec DI (Bonne pratique) :}
\begin{lstlisting}[language=CSharp, numbers=none]
public class MyController {
    private readonly StoreContext _context;
    public MyController(StoreContext context) { // Injecte!
        _context = context;
    }
}
\end{lstlisting}

\textbf{Avantages :}
\begin{itemize}
    \item \textbf{Testabilité} : On peut injecter un mock/fake pour les tests
    \item \textbf{Découplage} : Le controller ne sait pas comment créer le context
    \item \textbf{Lifecycle} : Le framework gère la durée de vie (Scoped, Singleton, Transient)
\end{itemize}
\end{defbox}

\begin{lstlisting}[language=CSharp, caption=GET - Récupérer les interventions avec filtrage par rôle]
[HttpGet]  // Repond aux requetes GET sur /api/demandeintervention
public async Task<ActionResult<PagedResult<object>>> GetDemandeInterventions(
    int page = 1, int pageSize = 5, string? search = null)
{
    // 1. Recuperer l'utilisateur depuis le cookie d'auth
    var currentUser = await _userManager.GetUserAsync(User);
    if (currentUser == null) return Unauthorized("User not found");

    // 2. Recuperer ses roles (un user peut avoir plusieurs roles)
    var roles = await _userManager.GetRolesAsync(currentUser);
    var role = roles.FirstOrDefault();  // Prend le premier

    // 3. Commencer la requete (pas encore executee!)
    var query = _context.DemandeInterventions.Where(d => d.IsActif);

    // 4. Appliquer le filtre selon le role
    if (role == "Demandeur")
    {
        query = query.Where(d => d.DemandeursID == userEntity.ID);
    }
    else if (role == "Utilisateur" || role == "Technicien")
    {
        query = query.Where(d => 
            d.DemandeursID == userEntity.ID || 
            d.DestinatairesID == userEntity.ID);
    }
    // Superviseur: pas de filtre = voit tout

    // 5. Eager Loading - charge les relations en 1 requete
    query = query
        .Include(d => d.Status)
        .Include(d => d.Demandeur)
        .Include(d => d.Destinataire);

    // 6. Pagination - SEULEMENT ICI la requete est executee
    var items = await query
        .Skip((page - 1) * pageSize)  // Saute les elements precedents
        .Take(pageSize)                // Prend N elements
        .ToListAsync();                // Execute la requete SQL

    return Ok(new PagedResult<object> { 
        Data = items, 
        TotalCount = totalCount 
    });
}
\end{lstlisting}

\begin{infobox}[title={\faLightbulb\ Concepts Avancés Expliqués}]
\textbf{1. Deferred Execution (Exécution Différée) :}
\begin{itemize}
    \item Les lignes avec \texttt{.Where()}, \texttt{.Include()} ne font PAS de requête SQL
    \item Elles construisent un \textbf{arbre d'expression} (Expression Tree)
    \item La requête SQL est générée et exécutée SEULEMENT au \texttt{ToListAsync()}
    \item Avantage : On peut composer la requête dynamiquement
\end{itemize}

\textbf{2. Le mot-clé async/await :}
\begin{itemize}
    \item \texttt{async} : Marque la méthode comme asynchrone
    \item \texttt{await} : "Attend" le résultat sans bloquer le thread
    \item \texttt{Task<T>} : Représente une opération asynchrone qui retournera T
    \item Permet de gérer des milliers de requêtes simultanées efficacement
\end{itemize}

\textbf{3. L'objet User :}
\begin{itemize}
    \item \texttt{User} est une propriété de \texttt{ControllerBase}
    \item Contient les \textbf{Claims} de l'utilisateur authentifié (depuis le cookie)
    \item \texttt{\_userManager.GetUserAsync(User)} récupère l'entité complète
\end{itemize}

\textbf{4. La pagination :}
\begin{itemize}
    \item \texttt{Skip((page - 1) * pageSize)} : Page 1 = Skip(0), Page 2 = Skip(5), etc.
    \item \texttt{Take(pageSize)} : Limite le nombre de résultats
    \item Génère : \texttt{OFFSET 0 ROWS FETCH NEXT 5 ROWS ONLY} en SQL
\end{itemize}
\end{infobox}

\newpage

\subsection{Service Angular : Intervention Service}

\begin{lstlisting}[language=TypeScript, caption=intervention.service.ts - Service Angular Complet]
@Injectable({
  providedIn: 'root'  // Singleton global - une seule instance
})
export class InterventionService {
  // URL de base construite depuis environment.ts
  private apiUrl = `${environment.apiUrl}/demandeintervention`;

  // Injection du HttpClient via le constructeur
  constructor(private http: HttpClient) {}

  // GET avec pagination et recherche
  getInterventions(
    page: number = 1,        // Valeur par defaut = 1
    pageSize: number = 10,   // Valeur par defaut = 10
    search?: string          // Optionnel (peut etre undefined)
  ): Observable<PagedResult<DemandeIntervention>> {
    
    // Construire les query parameters de maniere type-safe
    let params = new HttpParams()
      .set('page', page.toString())
      .set('pageSize', pageSize.toString());

    // Ajouter search seulement s'il existe
    if (search) {
      params = params.set('search', search);
    }

    // Faire la requete HTTP GET
    return this.http.get<PagedResult<DemandeIntervention>>(
      this.apiUrl,
      {
        params,              // Query params: ?page=1&pageSize=10
        withCredentials: true // Envoie le cookie d'auth
      }
    );
  }

  // POST avec upload de fichier optionnel
  createIntervention(
    intervention: Partial<InterventionDto>,  // Partial = tous champs optionnels
    file?: File                               // Fichier optionnel
  ): Observable<DemandeIntervention> {
    
    if (file) {
      // FormData permet d'envoyer fichier + JSON ensemble
      const formData = new FormData();
      formData.append('interventionData', JSON.stringify(intervention));
      formData.append('file', file);

      return this.http.post<DemandeIntervention>(
        this.apiUrl,
        formData,  // Content-Type: multipart/form-data (auto)
        { withCredentials: true }
      );
    }

    // Sans fichier: JSON simple
    return this.http.post<DemandeIntervention>(
      this.apiUrl,
      intervention,  // Content-Type: application/json (auto)
      { withCredentials: true }
    );
  }

  // Actions de workflow - POST sans body
  acceptIntervention(id: number): Observable<void> {
    return this.http.post<void>(
      `${this.apiUrl}/${id}/accept`,  // Template string pour l'URL
      {},                              // Body vide
      { withCredentials: true }
    );
  }
}
\end{lstlisting}

\begin{infobox}[title={\faLightbulb\ Concepts TypeScript/Angular Expliqués}]
\textbf{1. Décorateurs (@) :}
\begin{itemize}
    \item \texttt{@Injectable()} : Marque la classe comme injectable (peut recevoir des DI)
    \item \texttt{providedIn: 'root'} : Crée un \textbf{Singleton} au niveau de l'app
\end{itemize}

\textbf{2. Types Génériques (< >) :}
\begin{itemize}
    \item \texttt{Observable<PagedResult<DemandeIntervention>>} : Observable qui émet un PagedResult contenant des DemandeIntervention
    \item \texttt{http.get<T>()} : Le type de retour est typé comme T
\end{itemize}

\textbf{3. Partial<T> :}
\begin{itemize}
    \item \texttt{Partial<InterventionDto>} : Tous les champs de InterventionDto deviennent optionnels
    \item Utile pour les créations/mises à jour partielles
\end{itemize}

\textbf{4. Template Strings (`) :}
\begin{itemize}
    \item \texttt{\`\$\{this.apiUrl\}/\$\{id\}/accept\`} : Interpolation de variables dans une string
    \item Plus lisible que concaténation : \texttt{this.apiUrl + '/' + id + '/accept'}
\end{itemize}

\textbf{5. HttpParams (immutable) :}
\begin{itemize}
    \item \texttt{params.set()} retourne un \textbf{nouveau} objet (ne modifie pas l'original)
    \item C'est pourquoi on fait \texttt{params = params.set(...)}
\end{itemize}
\end{infobox}

\newpage

%==============================================================================
% SECTION 4 : MODULE CHATBOT IA
%==============================================================================
\section{\faRobot\ Module Chatbot IA - GMAO Chat}

\subsection{Architecture Text-to-SQL}

\begin{figure}[H]
\centering
\begin{tikzpicture}[
    scale=0.85,
    transform shape,
    node distance=0.8cm,
    step/.style={rectangle, rounded corners=8pt, draw=#1, fill=#1!20,
                 minimum width=11cm, minimum height=1.2cm, align=center, font=\small},
    arrow/.style={-{Stealth[length=3mm]}, thick, color=darkgray}
]

\node[step=primaryblue] (q) {\faUser\ "Quelles sont les interventions en cours ?"};
\node[step=accentorange, below=of q] (prep) {\faCog\ Prétraitement (normalisation, dates françaises)};
\node[step=codepurple, below=of prep] (llm) {\faBrain\ LLM génère le SQL via prompt engineering};
\node[step=warningred, below=of llm] (valid) {\faShieldAlt\ Validation sécurité (SELECT only)};
\node[step=darkgray, below=of valid] (exec) {\faDatabase\ Exécution sur SQL Server};
\node[step=successbox, below=of exec] (resp) {\faComment\ Formatage réponse en français};

\foreach \i/\j in {q/prep, prep/llm, llm/valid, valid/exec, exec/resp} {
    \draw[arrow] (\i) -- (\j);
}

\end{tikzpicture}
\caption{Pipeline de Traitement Text-to-SQL}
\end{figure}

\subsection{Configuration de l'API Key et des Providers}

\begin{lstlisting}[language=Python, caption=Configuration des providers IA]
# Dictionnaire de configuration pour chaque provider
AI_PROVIDERS = {
    "openrouter": {
        "api_key": "sk-or-v1-...",  # Cle secrete (ne jamais commit!)
        "endpoint": "https://openrouter.ai/api/v1/chat/completions",
        "model": "google/gemma-3-27b-it:free"  # Modele gratuit
    },
    "openai": {
        # .get() retourne '' si la cle n'existe pas (evite KeyError)
        "api_key": app_settings.get('OpenAI', {}).get('ApiKey', ''),
        "endpoint": "https://api.openai.com/v1/chat/completions",
        "model": "gpt-4o-mini"
    }
}

# Selection automatique du provider avec fallback
PRIMARY_AI_PROVIDER = None      # Le provider principal
FALLBACK_PROVIDERS = []         # Liste des providers de secours

for provider in ["openrouter", "gemini", "openai", "deepseek"]:
    # Verifie si l'API key existe et n'est pas vide
    if AI_PROVIDERS[provider]["api_key"]:
        if PRIMARY_AI_PROVIDER is None:
            PRIMARY_AI_PROVIDER = provider  # Premier trouve = principal
        else:
            FALLBACK_PROVIDERS.append(provider)  # Autres = fallback
\end{lstlisting}

\begin{infobox}[title={\faLightbulb\ Concepts Python Expliqués}]
\textbf{1. Dictionnaires Imbriqués :}
\begin{itemize}
    \item \texttt{AI\_PROVIDERS["openrouter"]["api\_key"]} : Accès à une valeur dans un dict imbriqué
    \item Structure : \texttt{\{"cle": \{"sous\_cle": "valeur"\}\}}
\end{itemize}

\textbf{2. La méthode .get() :}
\begin{itemize}
    \item \texttt{dict.get('key', default)} : Retourne \texttt{default} si la clé n'existe pas
    \item \texttt{dict['key']} : Lève \texttt{KeyError} si la clé n'existe pas
    \item \texttt{app\_settings.get('OpenAI', \{\}).get('ApiKey', '')} : Chaîne sécurisée de .get()
\end{itemize}

\textbf{3. Pattern Fallback :}
\begin{itemize}
    \item Si OpenRouter échoue, on essaie Gemini, puis OpenAI, puis DeepSeek
    \item Assure la \textbf{haute disponibilité} du service
    \item \texttt{is None} : Vérifie si la variable est exactement None (pas juste falsy)
\end{itemize}
\end{infobox}

\begin{warningbox}[title={\faKey\ Sécurité des API Keys}]
Les API Keys ne doivent \textbf{JAMAIS} être commitées sur Git. En production :
\begin{itemize}
    \item Utiliser des variables d'environnement
    \item Utiliser Azure Key Vault ou AWS Secrets Manager
    \item Configurer dans \texttt{appsettings.Production.json} (non versionné)
\end{itemize}
\end{warningbox}

\subsection{Prompt Engineering pour Text-to-SQL}

\begin{lstlisting}[language=Python, caption=Prompt système pour génération SQL]
def get_sql_generation_prompt() -> str:
    return f"""Tu es un expert SQL Server pour une application GMAO.
La date actuelle est: {datetime.now().strftime("%Y-%m-%d")}

## Schema de la base de donnees:
{DATABASE_SCHEMA}

## Regles:
1. Genere UNIQUEMENT des requetes SELECT
2. Utilise TOUJOURS des alias clairs
3. Pour les dates:
   - Aujourd'hui: CAST(GETDATE() AS DATE)
   - Comparaison: CAST(column AS DATE) = 'YYYY-MM-DD'
4. Pour les enregistrements actifs:
   - WHERE IsActif = 1 ou WHERE IsDeleted = 0

## Joins importants:
- DemandeInterventions -> Statuses ON StatusId = Id
- DemandeInterventions -> Users ON DemandeursID = ID
- Equipments -> EtatEquipements ON IdEtat = IdEtat

## Format:
Reponds UNIQUEMENT avec la requete SQL, sans markdown.
"""
\end{lstlisting}

\subsection{Validation de Sécurité SQL}

\begin{lstlisting}[language=Python, caption=Validation pour prévenir les injections SQL]
# Liste des mots-cles SQL dangereux (peuvent modifier/supprimer des donnees)
FORBIDDEN_KEYWORDS = [
    'INSERT', 'UPDATE', 'DELETE',  # Modification de donnees
    'DROP', 'TRUNCATE', 'ALTER', 'CREATE',  # Modification de schema
    'EXEC', 'EXECUTE',  # Execution de code
    'xp_', 'sp_',  # Procedures stockees systeme SQL Server
    '--', ';',     # Commentaires et multi-statements (injection classique)
    '/*', '*/'     # Commentaires bloc
]

def validate_sql(sql: str) -> bool:
    """Valide que le SQL est securise (SELECT uniquement)
    
    Args:
        sql: La requete SQL a valider
        
    Returns:
        bool: True si securise, False sinon
    """
    sql_upper = sql.upper()  # Comparaison insensible a la casse
    
    # Parcourir chaque mot-cle interdit
    for keyword in FORBIDDEN_KEYWORDS:
        if keyword in sql_upper:
            logger.warning(f"SQL bloque: mot-cle interdit '{keyword}'")
            return False  # Rejete immediatement
    
    # Verifier que ca commence par SELECT
    if not sql.strip().upper().startswith('SELECT'):
        logger.warning("SQL bloque: ne commence pas par SELECT")
        return False
    
    return True  # SQL valide
\end{lstlisting}

\begin{infobox}[title={\faLightbulb\ Concepts Python et Sécurité}]
\textbf{1. Type Hints (Annotations) :}
\begin{itemize}
    \item \texttt{def validate\_sql(sql: str) -> bool:} : Déclare les types attendus
    \item \texttt{sql: str} : Le paramètre doit être une string
    \item \texttt{-> bool} : La fonction retourne un booléen
    \item Aide l'IDE et les outils de vérification (mypy)
\end{itemize}

\textbf{2. Docstrings (\"\"\" \"\"\")} :
\begin{itemize}
    \item Documentation intégrée dans le code
    \item Accessible via \texttt{help(validate\_sql)} ou \texttt{validate\_sql.\_\_doc\_\_}
    \item Format Google/NumPy standard
\end{itemize}

\textbf{3. F-strings (f"...") :}
\begin{itemize}
    \item \texttt{f"SQL bloque: mot-cle interdit '\{keyword\}'"} : Interpolation de variables
    \item Plus lisible que \texttt{"SQL bloque: " + keyword}
    \item Peut contenir des expressions : \texttt{f"Total: \{a + b\}"}
\end{itemize}

\textbf{4. Protection contre les Injections SQL :}
\begin{itemize}
    \item \textbf{Attaque classique} : \texttt{'; DROP TABLE Users; --}
    \item On bloque \texttt{;} et \texttt{--} pour empêcher les multi-statements
    \item On force \texttt{SELECT} uniquement : lecture seule
\end{itemize}
\end{infobox}

\subsection{Flux de Communication Complet}

\begin{figure}[H]
\centering
\begin{tikzpicture}[
    scale=0.75,
    transform shape,
    node distance=1.5cm,
    comp/.style={rectangle, rounded corners=5pt, draw=#1, fill=#1!20,
                 minimum width=2.5cm, minimum height=1.2cm, align=center, font=\footnotesize},
    arrow/.style={-{Stealth[length=2.5mm]}, thick}
]

% Components
\node[comp=angularred] (angular) {\faAngular\\ Angular\\Chat Component};
\node[comp=dotnetpurple, right=2cm of angular] (dotnet) {\faServer\\ .NET API\\GMAOChatController};
\node[comp=accentorange, right=2cm of dotnet] (python) {\faPython\\ Python\\FastAPI Service};
\node[comp=codepurple, right=2cm of python] (llm) {\faBrain\\ OpenRouter\\LLM API};
\node[comp=darkgray, below=1.5cm of python] (db) {\faDatabase\\ SQL Server};

% Arrows with labels
\draw[arrow, angularred] (angular) -- node[above, font=\tiny] {1. POST /query} (dotnet);
\draw[arrow, dotnetpurple] (dotnet) -- node[above, font=\tiny] {2. Forward} (python);
\draw[arrow, accentorange] (python) -- node[above, font=\tiny] {3. Generate SQL} (llm);
\draw[arrow, codepurple] (llm) -- node[above, font=\tiny] {4. SQL} (python);
\draw[arrow, accentorange] (python) -- node[right, font=\tiny] {5. Execute} (db);
\draw[arrow, darkgray] (db) -- node[left, font=\tiny] {6. Results} (python);
\draw[arrow, accentorange, dashed] (python) -- node[below, font=\tiny] {7. Response} (dotnet);
\draw[arrow, dotnetpurple, dashed] (dotnet) -- node[below, font=\tiny] {8. JSON} (angular);

\end{tikzpicture}
\caption{Flux de Communication du Chatbot}
\end{figure}

\newpage

%==============================================================================
% SECTION 5 : AUTHENTIFICATION COMPLÈTE
%==============================================================================
\section{\faLock\ Système d'Authentification}

\subsection{Flux d'Authentification}

\begin{figure}[H]
\centering
\begin{tikzpicture}[
    node distance=0.6cm,
    step/.style={rectangle, rounded corners=5pt, draw=primaryblue, fill=primaryblue!15,
                 minimum width=12cm, minimum height=0.9cm, align=center, font=\small},
    num/.style={circle, draw=primaryblue, fill=primaryblue, text=white, 
                minimum size=0.5cm, font=\bfseries\footnotesize}
]

\node[step] (s1) {Utilisateur entre email + mot de passe sur la page Login};
\node[num, left=0.2cm of s1] {1};

\node[step, below=of s1] (s2) {Angular envoie POST /api/account/login avec credentials};
\node[num, left=0.2cm of s2] {2};

\node[step, below=of s2] (s3) {.NET vérifie via SignInManager.PasswordSignInAsync()};
\node[num, left=0.2cm of s3] {3};

\node[step, below=of s3] (s4) {Si OK : Cookie d'authentification créé automatiquement};
\node[num, left=0.2cm of s4] {4};

\node[step, below=of s4] (s5) {Retour JSON avec infos user (id, email, role)};
\node[num, left=0.2cm of s5] {5};

\node[step, below=of s5] (s6) {Angular stocke dans BehaviorSubject + localStorage};
\node[num, left=0.2cm of s6] {6};

\node[step, below=of s6] (s7) {Requêtes suivantes : Cookie envoyé automatiquement};
\node[num, left=0.2cm of s7] {7};

\foreach \i/\j in {s1/s2, s2/s3, s3/s4, s4/s5, s5/s6, s6/s7} {
    \draw[-{Stealth}, thick, primaryblue!70] (\i) -- (\j);
}

\end{tikzpicture}
\caption{Flux d'Authentification Complet}
\end{figure}

\subsection{Backend : AccountController}

\begin{lstlisting}[language=CSharp, caption=Login avec Cookie Authentication]
[HttpPost("login")]
public async Task<ActionResult> Login(LoginDto loginDto)
{
    // 1. Trouver l'utilisateur par email
    var user = await _userManager.FindByEmailAsync(loginDto.Email);
    if (user == null) return Unauthorized("Email incorrect.");

    // 2. Verifier le mot de passe
    var result = await _signInManager.PasswordSignInAsync(
        user.UserName, 
        loginDto.Password, 
        isPersistent: true,  // Cookie persistant
        lockoutOnFailure: false
    );
    
    if (!result.Succeeded) 
        return Unauthorized("Mot de passe incorrect.");

    // 3. Recuperer le role
    var roles = await _userManager.GetRolesAsync(user);
    var role = roles.FirstOrDefault();

    // 4. Recuperer l'entite User complete
    var userEntity = await _storeContext.Users
        .Include(u => u.Role)
        .FirstOrDefaultAsync(u => u.AppUserId == user.Id);

    // 5. Retourner les infos (le cookie est cree automatiquement)
    return Ok(new
    {
        userId = user.Id,
        email = user.Email,
        firstName = user.FirstName,
        lastName = user.LastName,
        role = role,
        internalId = userEntity.ID
    });
}
\end{lstlisting}

\subsection{Frontend : AuthService}

\begin{lstlisting}[language=TypeScript, caption=AuthService avec BehaviorSubject - Explique]
@Injectable({ providedIn: 'root' })  // Singleton global
export class AuthService {
  private apiUrl = `${environment.apiUrl}/account`;
  
  // BehaviorSubject: Observable avec valeur initiale qui garde la derniere valeur
  private currentUserSubject = new BehaviorSubject<UserInfo | null>(null);
  
  // Expose l'Observable en lecture seule (pas de .next() de l'exterieur)
  public currentUser$ = this.currentUserSubject.asObservable();

  constructor(private http: HttpClient) {
    this.checkAuthStatus();  // Verifie l'auth au demarrage de l'app
  }

  login(credentials: LoginDto): Observable<UserInfo> {
    return this.http.post<UserInfo>(
      `${this.apiUrl}/login`,
      credentials,
      { withCredentials: true }  // Important: envoie/recoit les cookies
    ).pipe(
      // tap: execute un effet secondaire sans modifier la valeur
      tap(user => {
        this.currentUserSubject.next(user);  // Met a jour le state
        localStorage.setItem('currentUser', JSON.stringify(user));
      })
    );
  }

  logout(): Observable<void> {
    return this.http.post<void>(
      `${this.apiUrl}/logout`,
      {},
      { withCredentials: true }
    ).pipe(
      tap(() => {
        this.currentUserSubject.next(null);  // Reset l'user
        localStorage.removeItem('currentUser');
      }),
      // catchError: gere les erreurs dans le pipe
      catchError(() => {
        this.currentUserSubject.next(null);
        localStorage.removeItem('currentUser');
        return of(void 0);  // of() cree un Observable qui emet une valeur
      })
    );
  }

  checkAuthStatus(): void {
    // 1. Restaurer depuis localStorage (instantane)
    const storedUser = localStorage.getItem('currentUser');
    if (storedUser) {
      this.currentUserSubject.next(JSON.parse(storedUser));
    }

    // 2. Verifier avec le serveur (asynchrone)
    this.http.get<UserInfo>(`${this.apiUrl}/user-info`, {
      withCredentials: true
    }).subscribe({
      next: (user) => this.currentUserSubject.next(user),
      error: () => {
        this.currentUserSubject.next(null);
        localStorage.removeItem('currentUser');
      }
    });
  }

  isAuthenticated(): boolean {
    return this.currentUserSubject.value !== null;
  }

  hasRole(role: string): boolean {
    return this.currentUserSubject.value?.role === role;
  }
}
\end{lstlisting}

\begin{infobox}[title={\faLightbulb\ Concepts RxJS et State Management}]
\textbf{1. BehaviorSubject vs Subject vs Observable :}
\begin{itemize}
    \item \texttt{Observable} : Flux de données en lecture seule
    \item \texttt{Subject} : Observable + peut émettre des valeurs (.next())
    \item \texttt{BehaviorSubject} : Subject + garde la dernière valeur + valeur initiale obligatoire
\end{itemize}

\textbf{2. Pattern Private Subject + Public Observable :}
\begin{lstlisting}[language=TypeScript, numbers=none]
private subject = new BehaviorSubject<T>(init);  // Prive: on peut .next()
public obs$ = this.subject.asObservable();       // Public: lecture seule
\end{lstlisting}
\begin{itemize}
    \item Empêche les composants d'appeler \texttt{.next()} directement
    \item Seul le service contrôle les mises à jour
\end{itemize}

\textbf{3. Opérateurs RxJS utilisés :}
\begin{itemize}
    \item \texttt{pipe()} : Chaîne des opérateurs (comme le pipe Unix |)
    \item \texttt{tap()} : Exécute un effet secondaire sans modifier la valeur
    \item \texttt{catchError()} : Intercepte les erreurs, peut retourner un Observable de remplacement
    \item \texttt{of()} : Crée un Observable qui émet une valeur puis se complète
\end{itemize}

\textbf{4. L'opérateur ?.} (Optional Chaining) :
\begin{itemize}
    \item \texttt{this.currentUserSubject.value?.role} : Retourne \texttt{undefined} si value est null
    \item Évite le crash : \texttt{Cannot read property 'role' of null}
\end{itemize}

\textbf{5. Pourquoi localStorage + BehaviorSubject ?}
\begin{itemize}
    \item \texttt{localStorage} : Persiste après refresh de la page
    \item \texttt{BehaviorSubject} : Réactivité en temps réel (les composants sont notifiés)
    \item Les deux ensemble : Meilleure UX
\end{itemize}
\end{infobox}

\newpage

%==============================================================================
% SECTION 6 : BASE DE DONNÉES
%==============================================================================
\section{\faDatabase\ Structure de la Base de Données}

\subsection{DbContext : StoreContext}

\begin{lstlisting}[language=CSharp, caption=Infrastructure/Data/StoreContext.cs]
public class StoreContext : IdentityDbContext<AppUser, IdentityRole, string>
{
    public StoreContext(DbContextOptions<StoreContext> options) 
        : base(options) { }
    
    // Tables principales
    public DbSet<Users> Users { get; set; }
    public DbSet<Role> Roles { get; set; }

    // Module Equipements
    public DbSet<Equipment> Equipments { get; set; }
    public DbSet<FamilleEquipment> FamilleEquipements { get; set; }
    public DbSet<EquipmentCriticite> EquipmentCriticites { get; set; }
    public DbSet<EtatEquipement> EtatEquipements { get; set; }

    // Module Interventions
    public DbSet<DemandeIntervention> DemandeInterventions { get; set; }
    public DbSet<TypeIntervention> TypeInterventions { get; set; }
    public DbSet<Urgence> Urgences { get; set; }
    public DbSet<Status> Statuses { get; set; }

    // Module Stock
    public DbSet<Stock> Stocks { get; set; }
    public DbSet<StockEntities> StockEntities { get; set; }
    public DbSet<StockStatus> StockStatuses { get; set; }
    public DbSet<StockSupplier> StockSuppliers { get; set; }

    protected override void OnModelCreating(ModelBuilder modelBuilder)
    {
        base.OnModelCreating(modelBuilder);
        
        // Appliquer les configurations Fluent API
        modelBuilder.ApplyConfiguration(new EquipmentConfiguration());
        modelBuilder.ApplyConfiguration(new DemandeInterventionConfiguration());
        // ... autres configurations
    }
}
\end{lstlisting}

\subsection{Diagramme Simplifié des Relations}

\begin{figure}[H]
\centering
\begin{tikzpicture}[
    scale=0.8,
    transform shape,
    table/.style={rectangle, draw=darkgray, fill=lightgray,
                  minimum width=3cm, minimum height=1.8cm, align=center, font=\small},
    relation/.style={-{Stealth[length=2mm]}, thick, darkgray}
]

% Tables
\node[table] (user) at (0, 0) {\textbf{Users}\\ID, Nom, Prenom\\Email, IDRole};
\node[table] (role) at (4, 0) {\textbf{Roles}\\IDRole, Libelle};
\node[table] (inter) at (0, -3) {\textbf{DemandeInterventions}\\ID, Titre, StatusId\\DemandeursID, DestinatairesID};
\node[table] (status) at (4, -3) {\textbf{Statuses}\\Id, Statut};
\node[table] (equip) at (8, 0) {\textbf{Equipments}\\ID, Designation\\IdEtat, IdCriticite};
\node[table] (etat) at (8, -3) {\textbf{EtatEquipements}\\IdEtat, NomEtat};

% Relations
\draw[relation] (user) -- node[above, font=\tiny] {N:1} (role);
\draw[relation] (inter) -- node[left, font=\tiny] {N:1} (user);
\draw[relation] (inter) -- node[above, font=\tiny] {N:1} (status);
\draw[relation] (equip) -- node[right, font=\tiny] {N:1} (etat);

\end{tikzpicture}
\caption{Relations Principales de la Base de Données}
\end{figure}

\newpage

%==============================================================================
% SECTION 7 : DEVOPS & FUTUR
%==============================================================================
\section{\faDocker\ DevOps et Améliorations Futures}

\subsection{Infrastructure Docker Actuelle}

\begin{lstlisting}[language=Python, caption=docker-compose.yml - Services]
services:
  # Base de donnees SQL Server
  sqlserver:
    image: mcr.microsoft.com/mssql/server:2022-latest
    environment:
      ACCEPT_EULA: "Y"
      SA_PASSWORD: "YourStrong!Password123"
    ports:
      - "1433:1433"
    volumes:
      - sqlserver_data:/var/opt/mssql
    healthcheck:
      test: /opt/mssql-tools18/bin/sqlcmd -S localhost -U sa ...
      interval: 10s
      retries: 10

  # Backend .NET
  backend:
    build:
      context: .
      dockerfile: API/Dockerfile
    environment:
      - ConnectionStrings__DefaultConnection=Server=sqlserver;...
      - ChatbotService__Url=http://chatbot:5002
    ports:
      - "5000:80"
    depends_on:
      sqlserver:
        condition: service_healthy

  # Chatbot Python
  chatbot:
    build: ./gmaochat
    ports:
      - "5002:5002"
    depends_on:
      - sqlserver

  # Frontend Angular
  frontend:
    build: ./client
    ports:
      - "4200:80"
    depends_on:
      - backend
\end{lstlisting}

\subsection{Pipeline CI/CD Proposé}

\begin{figure}[H]
\centering
\begin{tikzpicture}[
    scale=0.85,
    transform shape,
    node distance=0.5cm,
    stage/.style={rectangle, rounded corners=8pt, draw=#1, fill=#1!20,
                  minimum width=10cm, minimum height=1cm, align=center, font=\small},
    arrow/.style={-{Stealth[length=3mm]}, thick, darkgray}
]

\node[stage=primaryblue] (code) {\faCode\ Code Push sur GitHub/GitLab};
\node[stage=accentorange, below=of code] (build) {\faCog\ Build \& Tests Unitaires};
\node[stage=codepurple, below=of build] (docker) {\faDocker\ Build Images Docker};
\node[stage=darkgray, below=of docker] (push) {\faCloud\ Push vers Container Registry};
\node[stage=successbox, below=of push] (deploy) {\faServer\ Déploiement sur Azure/AWS};

\foreach \i/\j in {code/build, build/docker, docker/push, push/deploy} {
    \draw[arrow] (\i) -- (\j);
}

% Labels
\node[right=0.5cm of build, font=\footnotesize\color{gray}] {xUnit, Jest};
\node[right=0.5cm of docker, font=\footnotesize\color{gray}] {Multi-stage builds};
\node[right=0.5cm of push, font=\footnotesize\color{gray}] {ACR, ECR, Docker Hub};
\node[right=0.5cm of deploy, font=\footnotesize\color{gray}] {Kubernetes, App Service};

\end{tikzpicture}
\caption{Pipeline CI/CD Proposé}
\end{figure}

\subsection{Améliorations Futures Détaillées}

\begin{table}[H]
\centering
\rowcolors{2}{lightgray}{white}
\begin{tabular}{|>{\bfseries}l|l|l|}
\hline
\rowcolor{primaryblue}
\textcolor{white}{\textbf{Domaine}} & \textcolor{white}{\textbf{Amélioration}} & \textcolor{white}{\textbf{Outil}} \\
\hline
\multicolumn{3}{|c|}{\cellcolor{dockerblue!30}\textbf{CI/CD}} \\
\hline
Build & Automatisation des builds & GitHub Actions / GitLab CI \\
Tests & Tests automatisés & xUnit, Jest, Cypress \\
Qualité & Analyse de code & SonarQube \\
Déploiement & Déploiement continu & ArgoCD, Flux \\
\hline
\multicolumn{3}{|c|}{\cellcolor{successbox!30}\textbf{Infrastructure}} \\
\hline
Orchestration & Conteneurisation avancée & Kubernetes (AKS/EKS) \\
Scaling & Auto-scaling & HPA, Karpenter \\
Secrets & Gestion des secrets & Azure Key Vault \\
\hline
\multicolumn{3}{|c|}{\cellcolor{accentorange!30}\textbf{Monitoring}} \\
\hline
Logs & Centralisation logs & ELK Stack, Loki \\
Métriques & Métriques applicatives & Prometheus + Grafana \\
Tracing & Distributed tracing & Jaeger, Zipkin \\
Alerting & Alertes & PagerDuty, OpsGenie \\
\hline
\multicolumn{3}{|c|}{\cellcolor{warningred!30}\textbf{Sécurité}} \\
\hline
Scan & Scan de vulnérabilités & Trivy, Snyk \\
Secrets & Rotation des secrets & Vault \\
Réseau & Service Mesh & Istio, Linkerd \\
\hline
\end{tabular}
\caption{Roadmap DevOps}
\end{table}

\subsection{GitHub Actions - Exemple de Workflow}

\begin{lstlisting}[language=Python, caption=.github/workflows/ci-cd.yml (proposé)]
name: GMAO CI/CD Pipeline

on:
  push:
    branches: [main, develop]
  pull_request:
    branches: [main]

jobs:
  # Job 1: Build et Test Backend
  build-backend:
    runs-on: ubuntu-latest
    steps:
      - uses: actions/checkout@v4
      
      - name: Setup .NET
        uses: actions/setup-dotnet@v4
        with:
          dotnet-version: '9.0.x'
      
      - name: Restore dependencies
        run: dotnet restore
      
      - name: Build
        run: dotnet build --no-restore
      
      - name: Test
        run: dotnet test --no-build --verbosity normal

  # Job 2: Build et Test Frontend
  build-frontend:
    runs-on: ubuntu-latest
    steps:
      - uses: actions/checkout@v4
      
      - name: Setup Node.js
        uses: actions/setup-node@v4
        with:
          node-version: '20'
      
      - name: Install dependencies
        working-directory: ./client
        run: npm ci
      
      - name: Build
        working-directory: ./client
        run: npm run build --prod
      
      - name: Test
        working-directory: ./client
        run: npm test -- --watch=false --browsers=ChromeHeadless

  # Job 3: Build Docker Images
  docker:
    needs: [build-backend, build-frontend]
    runs-on: ubuntu-latest
    steps:
      - uses: actions/checkout@v4
      
      - name: Build and push Docker images
        run: |
          docker-compose build
          # docker-compose push (vers registry)
\end{lstlisting}

\newpage

%==============================================================================
% SECTION 8 : PATTERNS & BEST PRACTICES
%==============================================================================
\section{\faGraduationCap\ Patterns et Bonnes Pratiques}

\subsection{LINQ : Méthodes les Plus Utilisées}

\begin{table}[H]
\centering
\small
\rowcolors{2}{lightgray}{white}
\begin{tabular}{|>{\bfseries}p{3cm}|p{5cm}|p{5.5cm}|}
\hline
\rowcolor{dotnetpurple}
\textcolor{white}{\textbf{Méthode}} & \textcolor{white}{\textbf{Description}} & \textcolor{white}{\textbf{SQL Équivalent}} \\
\hline
Where(x => ...) & Filtre les éléments & WHERE condition \\
\hline
Select(x => ...) & Transforme/projette les éléments & SELECT colonnes \\
\hline
OrderBy(x => x.Prop) & Trie par ordre croissant & ORDER BY Prop ASC \\
\hline
OrderByDescending() & Trie par ordre décroissant & ORDER BY Prop DESC \\
\hline
Include(x => x.Nav) & Charge une relation & JOIN \\
\hline
FirstOrDefault() & Premier élément ou null & SELECT TOP 1 \\
\hline
Skip(n).Take(m) & Pagination & OFFSET n FETCH m \\
\hline
Count() & Compte les éléments & COUNT(*) \\
\hline
Any(x => ...) & Vérifie si au moins un existe & EXISTS \\
\hline
GroupBy(x => ...) & Groupe les éléments & GROUP BY \\
\hline
\end{tabular}
\caption{LINQ vers SQL - Correspondances}
\end{table}

\subsection{RxJS : Opérateurs Essentiels}

\begin{table}[H]
\centering
\small
\rowcolors{2}{lightgray}{white}
\begin{tabular}{|>{\bfseries}p{2.5cm}|p{5cm}|p{6cm}|}
\hline
\rowcolor{angularred}
\textcolor{white}{\textbf{Opérateur}} & \textcolor{white}{\textbf{Description}} & \textcolor{white}{\textbf{Exemple d'Usage}} \\
\hline
map() & Transforme chaque valeur & \texttt{.pipe(map(x => x * 2))} \\
\hline
filter() & Ne garde que certaines valeurs & \texttt{.pipe(filter(x => x > 0))} \\
\hline
tap() & Effet secondaire sans modifier & \texttt{.pipe(tap(x => console.log(x)))} \\
\hline
catchError() & Gère les erreurs & \texttt{.pipe(catchError(err => of(null)))} \\
\hline
switchMap() & Annule l'Observable précédent & Recherche avec debounce \\
\hline
mergeMap() & Fusionne les Observables & Requêtes parallèles \\
\hline
debounceTime() & Attend un délai d'inactivité & \texttt{.pipe(debounceTime(300))} \\
\hline
distinctUntilChanged() & Ignore les doublons consécutifs & Éviter requêtes inutiles \\
\hline
takeUntil() & Complète quand un autre émet & Unsubscribe automatique \\
\hline
\end{tabular}
\caption{Opérateurs RxJS les Plus Utiles}
\end{table}

\subsection{Pattern : Unsubscribe Automatique}

\begin{lstlisting}[language=TypeScript, caption=Éviter les memory leaks avec takeUntil]
@Component({ ... })
export class MyComponent implements OnInit, OnDestroy {
  // Subject qui emet quand le composant est detruit
  private destroy$ = new Subject<void>();

  ngOnInit(): void {
    // Toutes les subscriptions s'arretent automatiquement
    this.myService.getData().pipe(
      takeUntil(this.destroy$)  // Complete quand destroy$ emet
    ).subscribe(data => {
      this.data = data;
    });
  }

  ngOnDestroy(): void {
    this.destroy$.next();     // Emet une valeur
    this.destroy$.complete(); // Complete le Subject
  }
}
\end{lstlisting}

\begin{warningbox}[title={\faExclamationTriangle\ Memory Leaks en Angular}]
\textbf{Problème :} Si vous ne vous désabonnez pas des Observables, ils continuent à écouter même après la destruction du composant $\rightarrow$ \textbf{Memory Leak}.

\textbf{Solutions :}
\begin{itemize}
    \item \texttt{takeUntil(destroy\$)} : Pattern recommandé (ci-dessus)
    \item \texttt{async} pipe dans le template : Désabonnement automatique
    \item \texttt{subscription.unsubscribe()} : Manuel mais fastidieux
\end{itemize}
\end{warningbox}

\subsection{Pattern : Smart vs Dumb Components}

\begin{defbox}{Architecture de Composants}
\textbf{Smart Components (Containers) :}
\begin{itemize}
    \item Contiennent la logique métier
    \item Injectent les services
    \item Gèrent l'état
    \item Exemple : \texttt{InterventionsPageComponent}
\end{itemize}

\textbf{Dumb Components (Presentational) :}
\begin{itemize}
    \item Affichent des données via @Input()
    \item Émettent des événements via @Output()
    \item Pas de logique, pas de services
    \item Réutilisables et testables
    \item Exemple : \texttt{InterventionCardComponent}
\end{itemize}
\end{defbox}

\newpage

%==============================================================================
% SECTION 9 : GLOSSAIRE
%==============================================================================
\section{\faBookOpen\ Glossaire des Termes Techniques}

\begin{table}[H]
\centering
\small
\rowcolors{2}{lightgray}{white}
\begin{tabular}{|>{\bfseries}p{3.5cm}|p{10cm}|}
\hline
\rowcolor{primaryblue}
\textcolor{white}{\textbf{Terme}} & \textcolor{white}{\textbf{Définition}} \\
\hline
API & Application Programming Interface - Interface permettant la communication entre applications \\
\hline
REST & Representational State Transfer - Style d'architecture pour les APIs web (GET, POST, PUT, DELETE) \\
\hline
SPA & Single Page Application - Application web qui charge une seule page HTML et met à jour dynamiquement le contenu \\
\hline
ORM & Object-Relational Mapping - Technique pour convertir des données entre objets et bases de données relationnelles \\
\hline
DTO & Data Transfer Object - Objet utilisé pour transférer des données entre couches \\
\hline
DI & Dependency Injection - Pattern où les dépendances sont injectées plutôt que créées \\
\hline
Observable & Flux de données asynchrones pouvant émettre plusieurs valeurs dans le temps \\
\hline
BehaviorSubject & Type d'Observable qui garde la dernière valeur et l'émet aux nouveaux abonnés \\
\hline
Interceptor & Middleware qui intercepte les requêtes/réponses HTTP \\
\hline
Cookie Auth & Authentification basée sur des cookies HTTP sécurisés \\
\hline
CORS & Cross-Origin Resource Sharing - Mécanisme permettant les requêtes cross-domain \\
\hline
Text-to-SQL & Technique IA convertissant le langage naturel en requêtes SQL \\
\hline
Prompt Engineering & Art de concevoir des prompts efficaces pour les LLMs \\
\hline
CI/CD & Continuous Integration / Continuous Deployment - Automatisation du build et déploiement \\
\hline
Docker & Plateforme de conteneurisation d'applications \\
\hline
Kubernetes & Système d'orchestration de conteneurs \\
\hline
Lambda Expression & Fonction anonyme courte : \texttt{x => x.Property} \\
\hline
Async/Await & Syntaxe pour gérer le code asynchrone de manière synchrone \\
\hline
Middleware & Code qui s'exécute entre la requête et la réponse \\
\hline
\end{tabular}
\caption{Glossaire des Termes Techniques}
\end{table}

\newpage

%==============================================================================
% CONCLUSION
%==============================================================================
\section{\faFlagCheckered\ Conclusion}

\begin{tcolorbox}[
    enhanced,
    colback=primaryblue!10,
    colframe=primaryblue,
    title={\faStar\ Résumé de l'Application GMAO},
    fonttitle=\bfseries
]

\textbf{Architecture Technique :}
\begin{itemize}
    \item Backend \textbf{.NET Core 9} avec Entity Framework Core et Cookie Authentication
    \item Frontend \textbf{Angular 17+} avec architecture modulaire et RxJS
    \item Chatbot IA \textbf{Python FastAPI} avec approche Text-to-SQL
    \item Base de données \textbf{SQL Server} avec relations normalisées
\end{itemize}

\vspace{0.3cm}

\textbf{Modules Principaux :}
\begin{itemize}
    \item Gestion des \textbf{Interventions} avec workflow complet
    \item Gestion des \textbf{Équipements} et leur état
    \item Gestion des \textbf{Stocks} et fournisseurs
    \item \textbf{Chatbot IA} pour interroger les données en langage naturel
\end{itemize}

\vspace{0.3cm}

\textbf{Bonnes Pratiques Implémentées :}
\begin{itemize}
    \item Séparation des responsabilités (Controllers, Services, DTOs)
    \item Sécurité (Cookie Auth, validation SQL, filtrage par rôle)
    \item Architecture reactive (BehaviorSubject, Observables)
    \item Infrastructure Docker pour le déploiement
\end{itemize}

\end{tcolorbox}

\vspace{1cm}

\begin{center}
\begin{tikzpicture}
    \node[draw=successbox, fill=successbox!20, rounded corners=10pt,
          minimum width=12cm, minimum height=2cm, align=center] {
        {\Large\bfseries\color{successbox} Application GMAO Complète}\\[0.3cm]
        {\large Prête pour la démonstration devant le jury}\\[0.2cm]
        {\normalsize .NET Core $\bullet$ Angular $\bullet$ FastAPI $\bullet$ SQL Server $\bullet$ Docker}
    };
\end{tikzpicture}
\end{center}

%==============================================================================
% FIN DU DOCUMENT
%==============================================================================
\end{document}
